In this section I will present some of the current state-of-the-art solutions in binary-only grey-box fuzzing,
    distinguishing between \gls{SBI} and \gls{DBI} approaches.

    \begin{itemize}
        \item \underline{\textbf{\gls{DBI}}}
        \begin{itemize}
            \item \textbf{AFL++ in QEMU/FRIDA mode} \cite{aflplusplus}: AFL++ in QEMU mode uses QEMU in \textit{User-mode Emulation},
                which consists in translating guest machine code (i.e. guest \gls{BB}s) into an intermediate representation
                called \gls{TB} and subsequently appending the coverage tracing instructions at the end of it.
                This approach is based on emulating the entire CPU using QEMU's virtualized environment. As stated before,
                AFL in QEMU mode has overheads on the order of 600\% with respect to compile-time instrumentation \cite{zafl}.\\
                AFL++ in FRIDA mode \cite{aflplusplus-frida-mode}, instead, is based on shared library injection, in which a Frida \cite{frida} component
                called \textbf{Stalker} intercepts any \gls{BB} that is about to be executed, copies it into a \textbf{Code Cache}
                and dynamically instruments the copy with coverage tracing instructions.
                Execution control is then transferred to the copy, which is executed and then re-intercepted by \textbf{Stalker} at the end.
                Unlike QEMU mode, it does not emulate the CPU, but instead executes within the memory space of the target program.
                AFL++ reports that FRIDA mode is approximately 2 to 5 times slower than compile-time instrumentation \cite{aflplusplus_binary_only_targets}.

                \item \textbf{DynamoRIO/Pin} \cite{DynamoRIOPaper,PinPaper}: Both these approaches share the same basic principle, called
                \textbf{\gls{JIT} code caching}. This consists in taking control of the thread when the target application starts,
                and from that moment on, every \gls{BB} is intercepted before it is executed and is copied into a private \textbf{Code Cache}.
                The instrumentation is then added to the copy, and control is transferred to the code cache.
                The instrumentation process is more or less identical to FRIDA's Stalker mechanism, but is not as optimized.
                Also in this case, the implemented version uses AFL \cite{zalewski2017afl} as the fuzzer \cite{heuse2018afldynamorio,heuse2018aflpin}.
        \end{itemize}
            \item \underline{\textbf{\gls{SBI}}}
            \begin{itemize}
                \item \textbf{Retrowrite} \cite{retrowrite}: This tool performs symbolic lifting on the code, replacing every reference
                    discovered during disassembly with a symbolic label. Doing this allows instructions to be added without
                    modifying the references, achieving reassembleable assembly.
                    While RetroWrite can add ASan instrumentation, it cannot natively add AFL-style edge coverage instructions.
                    To do so, the lifted assembly must be instrumented with the afl-gcc tool in order to produce the patched binary.
                    Moreover, community-reported issues \cite{retrowrite_issue28} document missing instrumentation points when
                    RetroWrite's symbolized output is compiled through afl-gcc, caused by a mismatch between RetroWrite's
                    label-naming convention and afl-as's basic-block detection heuristics.
                    Finally, RetroWrite only works with PIE, non-stripped binaries.
                \item \textbf{AFL-Dyninst} \cite{afldyninst}: AFL-Dyninst employs Dyninst's binary rewriting capabilities
                        to instrument an existing executable by inserting callbacks at selected basic blocks.
                        Dyninst's instrumentation mechanism is based on code patching and trampoline redirection \cite{Dyninst},
                        which causes many context switches and state saving/restoring operations.
                        Its performance is slightly better than \gls{DBI}, but this approach still reaches 500\% fuzzing overhead,
                        as demonstrated by Stefan Nagy et al. \cite{zafl}.

                \item \textbf{ZAFL} \cite{zafl}: This tool uses two main phases: the first consists in
                    reconstructing an \gls{IR} using the Zipr binary rewriter.
                    The resulting \gls{IR} is then used by ZAX to perform an optimization on the \gls{IR}, find the instrumentation points,
                    and finally instrument these points. ZAFL works by inserting calls to a function that updates the shared map
                    at the beginning of each \gls{BB}.
                    Despite ZAFL's many improvements, Jifan Xiao et al. \cite{SystemCallPatternTracing} demonstrated that
                    it still causes crashes in 28\% of the binaries used in their cited paper.
            \end{itemize}
    \end{itemize}
    Before proceeding to the next section, in which I will discuss the main goals of this thesis, I want to discuss the core functioning of AFL and AFL++,
    given its importance for understanding my approach.
    \subsection{AFL Functioning}
    \cite{afl_technical_details}
        As previously stated, AFL and, in particular, AFL++ are among the most widely used grey-box
        fuzzing frameworks \cite{aflplusplus}.
        AFL++ represents the community-based continuation of the original AFL project.
        The core functioning is virtually the same and consists of a \gls{SHM} region shared between the fuzzer and the target.
        The \gls{SHM} size can vary, but the default is 65536 bytes.\\
        Each ``cell'' of the \gls{SHM} contains the number of times a given edge in the \gls{CFG} has been
            taken during the execution of the target. Since the \gls{SHM} has 65536 ``cells'', each ``cell''
            can contain a number from 0 to 255.
        The ``cell'' position is computed with a specific hash function applied to the IDs of the two \gls{BB}s
        that form the edge.
        The main Edge Tracing update rule is summarized in \cref{alg:afl_style_update}.
        \begin{algorithm}
            \caption{AFL-style SHM update rule}
            \label{alg:afl_style_update}
            \begin{algorithmic}[1]
                \STATE $idx \gets current\_bb\_id \oplus previous\_bb\_id$ \label{line:xor}
                \STATE $shared\_map[idx] \gets shared\_map[idx]+1$ \label{line:increment}
                \STATE $previous\_bb\_id \gets current\_bb\_id \gg 1$ \label{line:shift}
            \end{algorithmic}
        \end{algorithm}
        Line \cref{line:shift} of \cref{alg:afl_style_update} is used to avoid updating the same \gls{SHM}
        location when an edge is traversed forward or backward.
        An example of this could be a simple \texttt{do-while} loop, as seen in \cref{alg:do-while} and \cref{fig:Do-While-CFG}.\\
        \begin{figure}[ht]
            \centering
            \begin{minipage}[h]{0.48\textwidth}
                \centering
                \begin{algorithm}[H]
                    \caption{Do-While Example}
                    \label{alg:do-while}
                    \begin{algorithmic}
                        \REPEAT
                            \STATE \COMMENT{Basic Block A}
                        \UNTIL{$Permanence\_Condition$}
                        \STATE \COMMENT{Basic Block C}
                    \end{algorithmic}
                \end{algorithm}
            Without \cref{line:shift} of \cref{alg:afl_style_update}, the hash of Edge1
            would result in $ID\_A \oplus ID\_B$, and the hash of Edge2 (if the permanence condition is not met) would result
            in $ID\_B \oplus ID\_A = ID\_A \oplus ID\_B$, due to the symmetry of XOR. But since the hash is the same, there would be a hash collision
            that would cause the same ``cell'' to be updated twice instead of once.
            \end{minipage}
            \hfill
            \begin{minipage}[h]{0.48\textwidth}
                \centering
                \includegraphics[width=100pt]{AFLUpdateAlg3InstructionWhy.png}
                \caption{Do-while CFG}
                \label{fig:Do-While-CFG}
            \end{minipage}
        \end{figure}
        The main idea behind AFL is to instrument each basic block of a target application with the shared map update rule.
        The fuzzer initializes the \gls{SHM} memory with zeros, runs the instrumented target, which updates the \gls{SHM}
        according to \cref{alg:afl_style_update} until it reaches the end, and then returns to the fuzzer.
        The \gls{SHM} ID is communicated to the target through an environment variable called \textit{\_\_AFL\_SHM\_ID}.\\
        At this point the fuzzer checks the \gls{SHM} it has just obtained, called \texttt{trace\_bits}, searching for any newly discovered edge.
        This is done by keeping a structure called \texttt{virgin\_bits}, which contains all the edges not yet discovered and
        is initialized with ones. The comparison is performed through $\texttt{trace\_bits} \land \texttt{virgin\_bits}$. If any position results in one after the logical AND,
        this means a new edge has been discovered. In this case, the \texttt{virgin\_bits} data structure is updated and
        the input just executed is saved in the corpus. Otherwise, a new input is chosen, mutated and then executed. \\
        This approach requires a \texttt{fork()} call followed by an \texttt{execve()} call for each input,
        causing significant overhead for each input execution.
        Two main optimizations have been implemented:
        \begin{itemize}
            \item \textbf{forkserver}
            \item \textbf{persistent mode}: \small (available only in AFL++)
        \end{itemize}
        \vspace{0.5cm}
        \begin{center}
            \textbf{Forkserver}
        \end{center}
        The forkserver makes it possible to avoid calling \textbf{execve} for every input, performing it only once at the start of execution.
        The forkserver is instrumented before the main function of the target application.
        The AFL++ forkserver protocol can be summarized as follows:
        \begin{figure}[ht]
            \caption{AFL++ Forkserver Protocol}
            \centering
            \includegraphics[width=350pt]{AFLPlusPlusForkserverFunctioning.png}
            \label{fig:aflplusplusforkserver}
        \end{figure}
        \footnote{ \small \textit{The differences between the AFL forkserver protocol and the AFL++ forkserver protocol lie mainly in the messages exchanged between the forkserver and the fuzzer}}
        The communication between the Fuzzer and the Target shown in \cref{fig:aflplusplusforkserver} occurs
        unidirectionally through two different file descriptors created by the fuzzer:
        \begin{itemize}
            \item \textit{FORKSRV\_FD (fd 198)}: used by the forkserver to read a message sent by the fuzzer.
            \item \textit{FORKSRV\_FD + 1 (fd 199)}: used by the forkserver to write a message to the fuzzer.
        \end{itemize}
        \begin{enumerate}
            \item The forkserver sends the ``Welcome'' message to the fuzzer (0x41464c01).
            \item The fuzzer replies with the status code. At this point the forkserver must check whether $status \land 0xffffffff = 0x41464c01$.
                    If it is equal, execution continues; otherwise, the forkserver stops here.
            \item The forkserver sends the execution options to the fuzzer.
                These can indicate that the target is running in persistent mode, or other optimizations.
            \item The forkserver sends 0x41464c01 to the fuzzer to end the handshake.
            \item The fuzzer writes a message to the forkserver (its content is discarded); the handshake is now finished.
                        This message is the ``go'' signal for the forkserver.
            \item The forkserver invokes \texttt{fork()}. The parent writes the PID of the child to the fuzzer and then invokes \texttt{waitpid()}.\\
                    The child continues execution of the main function until it ends.
            \item When the parent returns from the \texttt{waitpid()} function (i.e. the child has terminated), the exit status
            of the child is sent to the fuzzer, and the forkserver waits for another ``go'' from the fuzzer.
        \end{enumerate}